\chapter{Sobre o Regulamento Geral}

\section{Vigência}
\begin{enumerate}
  \item O presente Regulamento encontra-se em vigor desde o momento em que aprovado e apresentado.
  \item O presente Regulamento é considerado caducado a partir do momento de entrada em vigor de Regulamento retificado.
\end{enumerate}

\section{Revisão}
\begin{enumerate}
  \item O presente Regulamento deve ser revisto sempre que considerado pertinente.
  \item Se retificado, o Regulamento Geral revisto deve ser proposto a aprovação.
  \item O processo de aprovação de uma retificação do Regulamento Geral deve reger-se pelos seguintes moldes:
  \begin{enumerate}
    \item Apresentação e discussão inicial em RDIR
    \item Aprovação por votação favorável de maioria da Direção da equipa.
    \item Apresentação da proposta de alteração à equipa em RGER, por um membro da Direção
    \item Período de reflexão de pelo menos 7 dias para todos os membros da equipa poderem analisar o documento fornecido
    \item Deverá haver um momento em RGER para esclarecimentos e discussão do novo regulamento antes da votação
    \item Após isto, o documento será aprovado em RGER caso obtenha voto favorável de pelo menos 75\% dos membros presentes. A versão retificada do Regulamento Geral entra em vigor imediatamente após a contagem dos votos ser confirmada.
    \item Caso rejeitado o documento pode ser revisto e apresentado para aprovação novamente, seguindo as diretrizes mencionadas nos pontos anteriores.
  \end{enumerate}
\end{enumerate}

\section{Casos Omissos}
\begin{enumerate}
  \item Qualquer eventualidade que não seja abrangida pelo presente Regulamento deve ser decidida pela Direção, segundo critérios de razoabilidade, proporcionalidade e funcionalidade.
  \item As decisões tomadas ao abrigo da presente secção devem ser registadas por escrito, incluindo a fundamentação da decisão e os membros intervenientes.
  \item Sempre que a decisão tenha impacto relevante no funcionamento da equipa, nos direitos ou deveres dos membros, ou na interpretação futura do presente Regulamento, deve ser posteriormente apresentada em RDIR ou RGER, conforme o seu âmbito, para conhecimento e eventual validação.
  \item A utilização recorrente desta secção para resolver situações semelhantes deve motivar a revisão do presente Regulamento, de modo a reduzir a existência de casos omissos.
\end{enumerate}

\section{Protocolos}
\begin{enumerate}
  \item Extensões ao Regulamento Geral podem ser executadas pela Board - denominadas de Protocolos.
  \item Os protocolos nunca devem isolada- ou coletivamente sobrepor-se ao presente Regulamento Geral, cuja autoridade é máxima.
\end{enumerate}
